% Lotka-Volterra phase portrait: closed orbits around the centre equilibrium.
% Vol VI Ch 8 §8.2.
\begin{figure}[htbp]
\centering
\begin{tikzpicture}
\begin{axis}[
  width=11cm,
  height=8cm,
  xlabel={Prey density $N$},
  ylabel={Predator density $P$},
  xmin=0, xmax=200,
  ymin=0, ymax=200,
  xtick={0, 50, 100, 150, 200},
  ytick={0, 50, 100, 150, 200},
  grid=major,
  grid style={draw=gray, very thin, opacity=0.3},
  every axis plot/.append style={very thick},
  samples=120,
]
  % Three closed orbits at different initial conditions (parametric ellipses)
  \addplot[engblue, smooth, no marks, domain=0:360]
    ({100 + 30*cos(x)}, {100 + 30*sin(x)});
  \addplot[engblue, smooth, no marks, domain=0:360]
    ({100 + 60*cos(x)}, {100 + 60*sin(x)});
  \addplot[engblue, smooth, no marks, domain=0:360]
    ({100 + 85*cos(x)}, {100 + 85*sin(x)});

  % Centre / equilibrium marker
  \addplot[engamber, mark=*, only marks, mark size=2.5pt]
    coordinates {(100, 100)};

  % Nullclines (horizontal predator-isocline, vertical prey-isocline)
  \addplot[engamber, dashed, no marks, domain=0:200] {100};
  \addplot[engamber, dashed, no marks] coordinates {(100, 0) (100, 200)};

  \node[anchor=west, engamber, font=\small] at (axis cs:108, 105)
    {$(N^*, P^*)$};
  \node[anchor=south, align=center, font=\footnotesize]
    at (axis cs:50, 175) {Closed orbits\\(neutral stability)};
\end{axis}
\end{tikzpicture}
\caption[Lotka-Volterra phase portrait]{Phase-portrait of the
Lotka-Volterra predator-prey system. The non-trivial equilibrium
$(N^*, P^*) = (\gamma/\delta, \alpha/\beta)$ is a centre rather
than a focus; trajectories are closed orbits whose amplitude is
set by the initial condition. The dashed lines are the prey and
predator nullclines; their intersection is the equilibrium.
Neutral stability is a structural prediction of the model that
real predator-prey systems do not display.}
\label{fig:vol06:ch08:lv-phase}
\end{figure}
